\section*{Bab \ref{ch:g11:polarity-and-cohesion} --- Keelektronegatifan, Kepolaran, dan Gaya Antarmolekul}

\begin{solution}{exo:g11:polarity-and-cohesion:1}
\ce{H-Cl}: klorin (3.16 dibandingkan 2.20). \ce{O-H}: oksigen (3.44).
\ce{C-O}: oksigen (3.44 dibandingkan 2.55).
\end{solution}

\begin{solution}{exo:g11:polarity-and-cohesion:2}
\ce{H-H}: 0, tidak polar. \ce{C-H}: 0.35, dianggap nonpolar. \ce{O-H}:
1.24, polar. \ce{N-H}: 0.84, polar. \ce{Cl-Cl}: 0, tidak polar.
\ce{C-Cl}: 0.61, polar.
\end{solution}

\begin{solution}{exo:g11:polarity-and-cohesion:3}
Nitrogen (3.04 > 2.19) dan oksigen (3.44 > 2.58): \omterm{def:g11:polarity-and-cohesion:electronegativity}{keelektronegatifan}
menurun ke bawah dalam \omterm{def:g9:periodic-table-first-look:period}{golongan}. Nitrogen (3.04 > 2.55) dan magnesium
(1.31 > 0.93): \omterm{def:g11:polarity-and-cohesion:electronegativity}{keelektronegatifan} meningkat dari kiri ke kanan
sepanjang periode.
\end{solution}

\begin{solution}{exo:g11:polarity-and-cohesion:4}
Amonia dan air: masing-masing memiliki \omterm{def:g7:atoms-and-molecules:atom}{atom} hidrogen yang terikat pada
nitrogen atau oksigen, dan \omterm{def:g10:lewis-and-shape:lone-pair}{pasangan elektron bebas} pada \omterm{def:g7:atoms-and-molecules:atom}{atom} itu. Metana
tidak memiliki hidrogen pada N, O, atau F; hidrogen klorida memiliki
hidrogen yang terikat pada klorin, yang tidak termasuk ketiga \omterm{def:g7:atoms-and-molecules:atom}{atom} dalam
definisi itu.
\end{solution}

\begin{solution}{exo:g11:polarity-and-cohesion:5}
Ya: \omterm{def:g11:polarity-and-cohesion:van-der-waals}{interaksi van der Waals} ada di antara semua \omterm{def:g7:atoms-and-molecules:molecule}{molekul}. Tetapi \omterm{def:g7:atoms-and-molecules:molecule}{molekul}
metana kecil (10 \omterm{def:g9:inside-the-atom:nucleus}{elektron}), sehingga interaksi itu sangat lemah dan
putus pada suhu yang sangat rendah: metana mendidih di sekitar
\qty{-162}{\celsius}.
\end{solution}

\begin{solution}{exo:g11:polarity-and-cohesion:6}
\ce{CH2Cl2}: polar (dua ikatan \ce{C-Cl} yang polar dan dua ikatan
\ce{C-H} yang hampir nonpolar: tidak ada simetri yang meniadakannya).
\ce{CCl4}: nonpolar (empat ikatan yang identik di sudut-sudut
tetrahedron). \ce{NH3}: polar (piramida). \ce{CO2}: nonpolar (linear,
ikatan yang berlawanan arah). \ce{BF3}: nonpolar (tiga \omterm{def:g11:polarity-and-cohesion:polar-bond}{ikatan polar}
yang identik dengan sudut \ang{120} dalam satu bidang saling
meniadakan).
\end{solution}

\begin{solution}{exo:g11:polarity-and-cohesion:7}
Hidrogen gugus \ce{O-H} dapat diberikan (bukan hidrogen pada \ce{CH3},
yang terikat pada karbon); \omterm{def:g7:atoms-and-molecules:atom}{atom} oksigen, dengan kedua \omterm{def:g10:lewis-and-shape:lone-pair}{pasangan elektron
bebasnya}, dapat menerima. Ya: \ce{O-H} pada metanol dapat berikatan
dengan oksigen sebuah \omterm{def:g7:atoms-and-molecules:molecule}{molekul} air, dan \ce{O-H} pada air dengan oksigen
metanol, misalnya \ce{CH3-O-H}\,$\cdots$\,\ce{OH2}.
\end{solution}

\begin{solution}{exo:g11:polarity-and-cohesion:8}
Ketiga \omterm{def:g7:atoms-and-molecules:molecule}{molekul} itu nonpolar dan sama bentuknya, tetapi makin lama makin
besar (10, 18, dan 36 \omterm{def:g9:inside-the-atom:nucleus}{elektron}): \omterm{def:g11:polarity-and-cohesion:van-der-waals}{interaksi van der Waals-nya} makin kuat,
demikian pula suhu didihnya.
\end{solution}

\begin{solution}{exo:g11:polarity-and-cohesion:9}
\omterm{def:g9:periodic-table-first-look:period}{Golongan} 14: metana tidak memiliki hidrogen yang terikat pada N, O,
atau F dan tidak membentuk \omterm{def:g11:polarity-and-cohesion:hydrogen-bond}{ikatan hidrogen}, sehingga metana mengikuti
kecenderungan \omterm{def:g9:periodic-table-first-look:period}{golongannya}.
\end{solution}

\begin{solution}{exo:g11:polarity-and-cohesion:10}
\omterm{def:g11:polarity-and-cohesion:van-der-waals}{Interaksi van der Waals}: \omterm{def:g7:atoms-and-molecules:molecule}{molekulnya} membesar dari \ce{HCl} ke \ce{HI},
dan peningkatan interaksi itu seiring ukuran lebih besar daripada
penurunan tarikan antara muatan-muatan parsialnya.
\end{solution}

\begin{solution}{exo:g11:polarity-and-cohesion:11}
\omterm{def:g7:atoms-and-molecules:molecule}{Molekul} iodin jauh lebih besar daripada \omterm{def:g7:atoms-and-molecules:molecule}{molekul} klorin (106 \omterm{def:g9:inside-the-atom:nucleus}{elektron}
dibandingkan 34): \omterm{def:g11:polarity-and-cohesion:van-der-waals}{interaksi van der Waals-nya} jauh lebih kuat, cukup
kuat untuk mengikat \omterm{def:g7:atoms-and-molecules:molecule}{molekul-molekulnya} dalam padatan pada suhu kamar.
\end{solution}

\begin{solution}{exo:g11:polarity-and-cohesion:12}
Propana nonpolar: hanya \omterm{def:g11:polarity-and-cohesion:van-der-waals}{interaksi van der Waals}. Metoksimetana polar
(dua ikatan \ce{C-O} yang polar dalam \ce{C-O-C} yang bengkok) tetapi
tidak memiliki hidrogen pada oksigennya: tidak ada \omterm{def:g11:polarity-and-cohesion:hydrogen-bond}{ikatan hidrogen} di
antara \omterm{def:g7:atoms-and-molecules:molecule}{molekul-molekulnya}. Etanol polar dan gugus \ce{O-H}-nya
membentuk \omterm{def:g11:polarity-and-cohesion:hydrogen-bond}{ikatan hidrogen}. Makin kuat tarikannya, makin tinggi suhu
didihnya: propana < metoksimetana < etanol.
\end{solution}

\begin{solution}{exo:g11:polarity-and-cohesion:13}
Jaringan es yang terbuka memerlukan ruang lebih besar daripada
\omterm{def:g7:atoms-and-molecules:molecule}{molekul-molekul} yang sama dalam cairan, sehingga massa jenis es lebih
kecil daripada air cair dan es terapung. Danau membeku dari atas:
lapisan es menyekat air di bawahnya, yang tetap cair, dan ikan-ikan
bertahan hidup.
\end{solution}

\begin{solution}{exo:g11:polarity-and-cohesion:14}
Metana padat (\omterm{def:g11:polarity-and-cohesion:van-der-waals}{interaksi van der Waals} antara \omterm{def:g7:atoms-and-molecules:molecule}{molekul-molekul} kecil) <
es (\omterm{def:g11:polarity-and-cohesion:hydrogen-bond}{ikatan hidrogen}) < kalium klorida (tarikan antara \omterm{def:g9:ions:ion}{ion-ion}).
\end{solution}

\begin{solution}{exo:g11:polarity-and-cohesion:15}
Sebuah \omterm{def:g7:atoms-and-molecules:molecule}{molekul} air memiliki dua \omterm{def:g7:atoms-and-molecules:atom}{atom} hidrogen untuk diberikan dan dua
\omterm{def:g10:lewis-and-shape:lone-pair}{pasangan elektron bebas} untuk menerima: setiap \omterm{def:g7:atoms-and-molecules:molecule}{molekul} dapat ikut serta
dalam empat \omterm{def:g11:polarity-and-cohesion:hydrogen-bond}{ikatan hidrogen}, dua diberikan dan dua diterima, yang
membangun jaringan yang lengkap. Sebuah \omterm{def:g7:atoms-and-molecules:molecule}{molekul} hidrogen fluorida
memiliki tiga \omterm{def:g10:lewis-and-shape:lone-pair}{pasangan elektron bebas} tetapi hanya satu \omterm{def:g7:atoms-and-molecules:atom}{atom} hidrogen:
\omterm{def:g7:atoms-and-molecules:molecule}{molekul} itu memberikan satu \omterm{def:g11:polarity-and-cohesion:hydrogen-bond}{ikatan hidrogen} saja, sehingga rata-rata
hanya satu ikatan per \omterm{def:g7:atoms-and-molecules:molecule}{molekul} yang dapat terbentuk (rantai zig-zag).
Air memiliki kira-kira dua kali lebih banyak \omterm{def:g11:polarity-and-cohesion:hydrogen-bond}{ikatan hidrogen} per
\omterm{def:g7:atoms-and-molecules:molecule}{molekul}, dan mendidih pada suhu yang lebih tinggi.
\end{solution}

\begin{solution}{pb:g11:polarity-and-cohesion:1}
\textbf{1.} \ce{H-O-H}, \ce{H-S-H}, \ce{H-Se-H}, setiap \omterm{def:g7:atoms-and-molecules:atom}{atom} pusat
dengan dua \omterm{def:g10:lewis-and-shape:lone-pair}{pasangan elektron bebas}: O, S, dan Se memiliki enam \omterm{def:g10:electron-shells:valence}{elektron
valensi}, karena berada dalam \omterm{def:g9:periodic-table-first-look:period}{golongan} yang sama.

\textbf{2.} $M(\ce{H2O}) = \qty{18.0}{g/mol}$, $M(\ce{H2S}) =
\qty{34.1}{g/mol}$, $M(\ce{H2Se}) = \qty{81.0}{g/mol}$.

\textbf{3.} $3.44 - 2.20 = 1.24$; $2.58 - 2.20 = 0.38$;
$2.55 - 2.20 = 0.35$. Hanya ikatan \ce{O-H} yang jelas polar.

\textbf{4.} Bengkok (empat kelompok, dua di antaranya \omterm{def:g10:lewis-and-shape:lone-pair}{pasangan bebas}).
Air polar; hidrogen sulfida dan hidrogen selenida hanya sangat sedikit.

\textbf{5.} 10, 18, dan 36 \omterm{def:g9:inside-the-atom:nucleus}{elektron}: hidrogen selenida memiliki
\omterm{def:g11:polarity-and-cohesion:van-der-waals}{interaksi van der Waals} yang paling kuat.

\textbf{6.} Hanya air: \omterm{def:g7:atoms-and-molecules:atom}{atom-atom} hidrogennya terikat pada oksigen, salah
satu dari ketiga \omterm{def:g7:atoms-and-molecules:atom}{atom} yang sangat elektronegatif; belerang dan selenium
tidak cukup elektronegatif untuk memberi \omterm{def:g7:atoms-and-molecules:atom}{atom-atom} hidrogennya
$\delta^+$ yang nyata.

\textbf{7.} Empat: dua melalui \omterm{def:g7:atoms-and-molecules:atom}{atom-atom} hidrogennya sendiri, dua
diterima oleh kedua \omterm{def:g10:lewis-and-shape:lone-pair}{pasangan elektron bebasnya}.

\textbf{8.} \omterm{def:g11:polarity-and-cohesion:van-der-waals}{Interaksi van der Waals} (termasuk tarikan lemah antara
muatan-muatan parsial yang kecil).

\textbf{9.} $212.9 - 273.15 = \qty{-60.3}{\celsius}$ dan
\qty{-41.3}{\celsius}.

\textbf{10.} $-41.3 - (-60.3) = \qty{19.0}{\celsius}$.

\textbf{11.} \omterm{def:g7:atoms-and-molecules:molecule}{Molekulnya} lebih besar, sehingga \omterm{def:g11:polarity-and-cohesion:van-der-waals}{interaksi van der Waals-nya}
lebih kuat.

\textbf{12.} $-60.3 - 19.0 = \qty{-79.3}{\celsius}$.

\textbf{13.} Kenaikan $-88.1 - (-112) = \qty{23.9}{\celsius}$; ramalan
untuk metana $-112 - 23.9 = \qty{-136}{\celsius}$, dibandingkan
\qty{-162}{\celsius} yang sebenarnya: ekstrapolasinya sekitar
\qty{26}{\celsius} terlalu tinggi. Ekstrapolasi itu tidak tepat;
kecenderungan yang sebenarnya bukan garis lurus.

\textbf{14.} Kenaikan $\qty{18.7}{\celsius}$; ramalan untuk hidrogen
fluorida $-85.1 - 18.7 = \qty{-103.8}{\celsius}$, dibandingkan
\qty{19.6}{\celsius}: selisih sekitar \qty{123}{\celsius}.

\textbf{15.} $100 - (-79.3) = \qty{179}{\celsius}$, sekitar
\qty{180}{\celsius}.

\textbf{16.} Kenaikan $\qty{25.3}{\celsius}$; ramalan $-87.8 - 25.3 =
\qty{-113}{\celsius}$; selisih $-33.4 - (-113) = \qty{80}{\celsius}$.

\textbf{17.} Air (\qty{180}{\celsius}) > hidrogen fluorida
(\qty{123}{\celsius}) > amonia (\qty{80}{\celsius}). Air memberikan dua
\omterm{def:g7:atoms-and-molecules:atom}{atom} hidrogen dan menerima dengan dua \omterm{def:g10:lewis-and-shape:lone-pair}{pasangan elektron bebas}: empat
\omterm{def:g11:polarity-and-cohesion:hydrogen-bond}{ikatan hidrogen} per \omterm{def:g7:atoms-and-molecules:molecule}{molekul}, semuanya terpakai. Hidrogen fluorida hanya
memiliki satu hidrogen untuk diberikan, dan amonia hanya satu \omterm{def:g10:lewis-and-shape:lone-pair}{pasangan
elektron bebas} untuk menerima: rata-rata lebih sedikit ikatan per
\omterm{def:g7:atoms-and-molecules:molecule}{molekul}, dan nitrogen, yang kurang elektronegatif, membentuk ikatan yang
lebih lemah.

\textbf{18.} Gas: air akan mendidih di sekitar \qty{-80}{\celsius}. Bumi
tidak akan memiliki air cair, tidak ada samudra, tidak ada hujan, dan
tidak ada kehidupan yang bergantung padanya.

\textbf{19.} Jawaban itu mengekstrapolasi garis lurus yang ditarik
melalui dua titik saja, dan uji pada \omterm{def:g9:periodic-table-first-look:period}{golongan} 14 menunjukkan bahwa
ekstrapolasi seperti itu dapat meleset sekitar \qty{25}{\celsius}.

\textbf{20.} Tanpa \omterm{def:g11:polarity-and-cohesion:hydrogen-bond}{ikatan hidrogen}, air akan mendidih pada sekitar
\qty{-80}{\celsius}, kira-kira \qty{180}{\celsius} di bawah titik
didihnya yang sebenarnya.
\end{solution}
